\documentclass[9pt]{extarticle}

% to be able to add references and citations and links in the document
\usepackage[colorlinks=true, linkcolor=blue, urlcolor=cyan]{hyperref}

% to set custom margins
\usepackage[left = 1cm, right = 1cm, bottom = 1cm, top = 0.5cm]{geometry}

% to be able to use color in document
\usepackage[dvipsnames]{xcolor}
\definecolor{headings}{HTML}{701112} % Dark red color for headings

% to add icons
\usepackage{fontawesome5}
\newcommand{\seticon}[1]{\textcolor{black}{\csname #1\endcsname}}

% to have old-style figures
% \usepackage[osf]{ebgaramond}

% for adding footer (date)
\usepackage{fancyhdr}

% to create colourful boxes
\usepackage[skins,many]{tcolorbox}

% to write elegant math
\usepackage{amsmath}

% to write math symbols
\usepackage{amssymb}

% to be able to quote things verbatim
\usepackage{listings}

% to be able to place table at exact location
\usepackage{float}

% to be able to use top rule etc in tables
\usepackage{booktabs}

% to have no space in all lists
\usepackage{enumitem}
\setlist[enumerate]{nosep}
\setlist[itemize]{nosep}

% to remove section numbering
\setcounter{secnumdepth}{0}

% to modify how titles appear
\usepackage{titlesec}
% remove title spacing above and below
\titlespacing{\section}{0pt}{*0}{*0}
\titlespacing{\subsection}{0pt}{*0}{*0}

% to not have numbering
\pagestyle{empty}

\begin{document}
    %% HEADER %%
    \begin{center}
        \Huge
        D\textsc{anish} S\textsc{hakeel} M\textsc{ohammad} \\
    \end{center}
    % \vspace{-14pt}
    \begin{center}
        % \small
        \seticon{faEnvelope} mohammadsdanish@gmail.com \hspace{0.5em} $\vert$ \hspace{0.5em}
        \seticon{faPhone} (+44) 078237-81284 \hspace{0.5em} $\vert$ \hspace{0.5em}
        \seticon{faGithub} \href{https://github.com/TheDansmod}{github.com/TheDansmod} \\ % $\vert$
        \seticon{faLinkedin} \href{https://www.linkedin.com/in/danish-mohammad/}{www.linkedin.com/in/danish-mohammad/} \hspace{0.5em} $\vert$ \hspace{0.5em}
        \seticon{faHome} Zone 2, East London, E2 9HU \\
        Full UK work authorization (Graduate visa) until April 2028, no sponsorship required.\\
        \vspace{0.5em}
        \noindent\fbox{%
            \parbox{0.87\textwidth}{%
                Python developer with four years of applied development experience spanning production software engineering, high-performance simulation systems, and statistical modelling. Demonstrated capacity to design and implement large-scale search and simulation engines (20M+ state transitions/second), rebuild production data pipelines, and apply rigorous benchmarking methodology to novel algorithms. Seeking to apply this combination of systems-engineering rigor and quantitative research to the development of production trading strategies and simulation engines.
            }%
        }
    \end{center}

    %% EDUCATION %%
    \vspace{1pt}
    \noindent {\large \textcolor{headings}{\textbf{EDUCATION}}} \\
    \vspace{-7pt}
    \hrule
    \vspace{5pt}
    \noindent {\normalsize \textbf{Queen Mary University of London}}, London, UK \hfill 2024 -- 2026 \\
    {\normalsize Master of Science in Artificial Intelligence $\vert$ Merit}

    \vspace{5pt}
    \noindent {\normalsize \textbf{Birla Institute of Technology and Science}}, Pilani, India \hfill 2018 -- 2022 \\
    {\normalsize Bachelor of Engineering in Computer Science $\vert$ First Division}

    %% PROJECTS %%
    \vspace{10pt}
    \noindent {\large \textcolor{headings}{\textbf{PROJECTS}}} \\
    \vspace{-7pt}
    \hrule

    \vspace{5pt}
    \noindent {\normalsize\textbf{High-Throughput Simulation Engine for Exhaustive Game Solving (Tinyhouse Chess)}} \hfill Sep 2026 - present
    \begin{itemize}[leftmargin=*]
        \item Exhaustive game-theoretic solution of a combinatorial game requires enumerating the full reachable state space via breadth-first search and subsequently propagating terminal-position values backward through the state graph (retrograde analysis); at full-game scale, the binding constraint is not move-generation compute but sustained write throughput to persistent storage.
        \item Designed and implemented a Python-based move-generation and state-transition engine sustaining 20 million plies/second, through bitboard state representation, batched I/O with a sharded log-structured merge architecture to manage the I/O bottleneck.
        \item Currently engineering a high-write-throughput persistence layer to support exhaustive state enumeration at full-game scale, as the prerequisite stage before implementing retrograde analysis to obtain a complete win/loss/draw classification over the entire reachable state space.
        \item Directly transferable to production simulation-engine design: both problem classes require correctness-verified, high-throughput state transitions over large discrete state spaces under tight memory and I/O constraints, with performance bounded by the slower one between compute throughput and sustained write throughput.
    \end{itemize}

    \vspace{5pt}
    \noindent {\normalsize \textbf{FRAG: A Multi-modal Multi-agent Fashion Recommendation RAG Agent} \href{https://github.com/TheDansmod/fashion_mas}{demo \faExternalLink*}} \hfill Dec 2025 -- Mar 2026
    \begin{itemize}[leftmargin=*]
        \item Built an end-to-end agentic Python system recommending items from a 300k-item catalogue using combined text and image inputs, addressing the inability of standard recommendation pipelines to jointly reason over visual and textual preference signals.
        \item Designed a LangGraph-orchestrated multi-agent architecture with a Self-RAG-inspired self-evaluation loop; implemented a custom MCP server backed by a Qdrant vector database, with a persistent, streaming conversational interface.
        \item Implemented production-style monitoring (LangSmith) to continuously track retrieval quality, generation failure modes, and agent-loop latency; used this telemetry to diagnose and close two systematic failure modes prior to final evaluation. This constitutes a monitoring/diagnosis workflow directly analogous to live strategy performance monitoring.
        \item Designed a 60-query adversarial evaluation dataset and a calibrated LLM-as-a-judge scoring rubric (relevance, coherence, diversity) to obtain a defensible, bias-resistant evaluation signal; achieved 9.3/10 under this framework — evidence of rigorous, quantitative evaluation-design methodology applicable to strategy back-testing.

    \end{itemize}

    % \vspace{5pt}
    % \noindent {\normalsize \textbf{MCTS Improvements for Sushi-Go}} \hfill Sep 2024 -- Nov 2024 (QMUL) \\
    % \textit{Course}: AI in Games
    % \begin{itemize}[leftmargin=*]
    %     \item Implementation and comprehensive experimental evaluation of the impact of \href{https://link.springer.com/article/10.1007/s10462-022-10228-y}{improvements} (like using novel heuristics, progressive bias, and RAVE) to the Monte Carlo Tree Search (MCTS) algorithm within the \href{https://gaigresearch.github.io/projects/TAG}{Tabletop Games Framework} (TAG) for the game of Sushi Go!
    %     \item Obtained a 12\% improvement over the existing advanced MCTS agent.
    % \end{itemize}

    %% WORK EXPERIENCE %%
    \vspace{10pt}
    \noindent {\large \textcolor{headings}{\textbf{WORK EXPERIENCE}}} \\
    \vspace{-7pt}
    \hrule
    \vspace{5pt}
    \noindent {\normalsize \textbf{Cisco Systems}, India \(\vert\) Software Engineer} \hfill 10 Aug 2022 -- 12 Sep 2024 (FT)
    \begin{itemize}[leftmargin=*]
        \item Software Engineer on Cisco's Identity Services Engine (ISE), a network access control and policy enforcement platform deployed across 10,000+ enterprises globally; work centred on production correctness and reliability of a large-scale asynchronous distributed system.
        \item Defined acceptance criteria and edge-case validation logic to determine which ML-generated device-classification proposals were reliable enough for production deployment, reducing erroneous auto-classification in previously unclassified device categories. This constitutes direct experience in setting production-readiness criteria for model-driven automated decisions.
        \item Conducted systematic failure analysis across 80+ production defects spanning database consistency anomalies, concurrent race conditions, and policy-evaluation edge cases; built reproducible test cases isolating root causes, converting ambiguous production failures into structured, tractable engineering problems — a discipline directly applicable to testing and hardening production trading strategies.
    \end{itemize}

    \vspace{5pt}
    \noindent {\normalsize \textbf{Epsilon Data Management}, India \(\vert\) Analyst Intern} \hfill 08 Jun 2021 -- 09 Aug 2021 (FT)
    \begin{itemize}[leftmargin=*]
        \item The data team relied on proprietary SAS code for customer segmentation followed by manual modification in Excel, which was slow and inaccessible to Python-based workflows. Tasked with rebuilding this pipeline in Python during a 2-month internship.
        % \item Designed and developed a scalable, Python-based translation of the SAS segmentation pipeline, ensuring functional parity and improved maintainability. Reduced pipeline processing time by 43\%, enabling faster iteration for the analytics team. - rewrite this for micro1
        \item Rebuilt the SAS customer segmentation pipeline in Python, ensuring functional parity through systematic output validation at each transformation stage; reduced pipeline processing time by 43\% while establishing a QA layer that enabled the analytics team to audit segmentation decisions with explicit, reproducible logic rather than opaque SAS macros.
    \end{itemize}

    % \vspace{5pt}
    % \noindent {\normalsize \textbf{Parentlane}, India \(\vert\) Machine Learning Intern} \hfill 18 May 2020 -- 28 Jun 2020 (FT)
    % \begin{itemize}[leftmargin=*]
    %     \item Leveraged a pre-existing question-answer and article corpus along with transformer based article summarisation and similarity based question matching to develop automated content suggestion and answer retrieval for user queries.
    % \end{itemize}

    %% AWARDS %% - listing just a single one is not beneficial - either integrate it into other locations or remove it or add more
    % \vspace{10pt}
    % \noindent {\large \textcolor{headings}{\textbf{AWARDS}}} \\
    % \vspace{-7pt}
    % \hrule
    % \vspace{5pt}
    % \noindent {\normalsize \textbf{Queen Mary Global Talent Scholarship} \hfill 2024 \\
    % Queen Mary University of London

    %% MASTER'S THESIS %%
    \vspace{10pt}
    \noindent {\large \textcolor{headings}{\textbf{MASTER'S THESIS}}} \\
    \vspace{-7pt}
    \hrule
    \vspace{5pt}
    \noindent {\normalsize \textbf{Causal Discovery in the presence of Latent Confounders}} \hfill May 2025 -- Dec 2025 (QMUL) \\
    \textit{Advisor}: Dr. Anthony Constantinou \hfill \href{https://github.com/TheDansmod/qmul-msc-ai-code-master-thesis/blob/main/dissertation_report.pdf}{paper \faExternalLink*} $\vert$ \href{https://github.com/TheDansmod/qmul-msc-ai-code-master-thesis}{code \faExternalLink*}
    \begin{itemize}[leftmargin=*]
        \item Addressed the problem of recovering cause–effect structure between variables when some variables are unobserved — a setting common in real-world, non-experimental data.
        \item Contributed two novel mapping functions converting a discrete combinatorial-optimisation problem into a continuous one, enabling the use of gradient-free black-box optimisers (PPO, CMA-ES) in place of combinatorial search.
        \item Designed a rigorous benchmarking protocol against state-of-the-art causal discovery algorithms across synthetic and real-world graphs of varying confounder density; validated that a 30\% reduction in computational cost did not degrade structural accuracy beyond acceptable margins, establishing the precise operating conditions under which the proposed optimisation is appropriate — a methodology directly analogous to validating a simulation or backtesting engine against a reference implementation.
    \end{itemize}

    %% BACHELOR'S THESIS %%
    \vspace{10pt}
    \noindent {\large \textcolor{headings}{\textbf{ADDITIONAL PROJECTS}}} \\
    \vspace{-7pt}
    \hrule
    \vspace{5pt}
    \noindent {\normalsize \textbf{Bachelor's Thesis: Application of Actor-Critic Methods to Improving Video Quality Experience}} \hfill Jan -- May 2022 (BITS) \\
    \textit{Advisor}: Dr. Paresh Saxena
    \begin{itemize}[leftmargin=*]
        \item Adaptive Bitrate (ABR) algorithms in streaming video players must balance bandwidth efficiency and user Quality of Experience (QoE).
        \item Investigated whether Deep RL actor-critic methods could outperform the then state-of-the-art model (\href{https://dl.acm.org/doi/10.1145/3098822.3098843}{Pensieve}) for server-side ABR.
        % \item Implemented AR-A3C. It proved unstable, so pivoted to PPO with optimisations like advantage normalisation and value function clipping. Outperformed Pensieve by approximately 30\% on QoE metrics, validating the approach.
        \item Initially implemented AR-A3C; upon diagnosing training instability as arising from high-variance policy gradients in the streaming reward landscape rather than hyperparameter misconfiguration, made the principled decision to pivot to PPO with advantage normalisation and value function clipping. Outperformed Pensieve by approximately 30\% on QoE metrics, validating the approach.
    \end{itemize}

    \vspace{5pt}
    \noindent {\normalsize \textbf{Sudoku Puzzle Generation with Deep Reinforcement Learning}} \hfill Feb 2025 -- Apr 2025 \\
    \textit{Course}: Computational Creativity (Prof. Simon Colton)
    \begin{itemize}[leftmargin=*]
        \item Applied deep RL (DDQN, Rainbow) with convolutional function approximators to autonomously generate valid Sudoku puzzles of controlled, increasing difficulty, addressing sparse-reward shaping via a potential-based reward function.
        \item Designed a quantitative difficulty-evaluation rubric and validated that generated puzzles achieved a difficulty score 67\% above a randomised baseline, confirming the policy had learned a coherent difficulty-shaping strategy rather than a degenerate heuristic.
    \end{itemize}

    %% TECHNICAL SKILLS %%
    \vspace{10pt}
    \noindent {\large \textcolor{headings}{\textbf{TECHNICAL SKILLS}}} \\
    \vspace{-7pt}
    \hrule
    \vspace{5pt}
    \noindent 
    \begin{tabular}{l p{15cm}}
        \textbf{Software} & Python (primary), high-performance/algorithmic programming, data structures, simulation engine design, Java, SQL, Git, Docker, Unix \\% , Java Spring, RabbitMQ, Redis\\
        \textbf{Quant \& Stat Methods} & reinforcement learning (PPO, DDQN, Rainbow), evolutionary strategies (CMA-ES), reward shaping, Monte Carlo-style simulation design, statistical evaluation and benchmarking design \\
        \textbf{Data Engr} & pipeline construction and validation, ETL, pandas, numpy, sklearn, SQLite, Qdrant (vector DB) \\
        \textbf{ML / AI} & PyTorch, LangChain/LangGraph, LangSmith, Chainlit, MCP (Model Context Protocol), Hydra, Ollama, jupyter \\
        \textbf{Other} & Git, Jenkins, Docker, Unix, \LaTeX \\
    \end{tabular}

    \pagestyle{fancy}
    \fancyhf{}
    \renewcommand{\footrulewidth}{0pt}
    \renewcommand{\headrulewidth}{0pt}
    % \rfoot{\normalsize Date: \textit{20 Mar, 2026}}
\end{document}
